\documentclass[11pt]{amsart}

\usepackage[T1]{fontenc}
\usepackage{lmodern}
\usepackage{microtype}
\usepackage{mathtools,amssymb,amsthm}
\usepackage{textcomp}
\usepackage{upquote}
\usepackage[hidelinks]{hyperref}

\raggedbottom

\hypersetup{
  pdftitle={A nine-vertex P4-sparse minimal (2,3)-polar obstruction that is not a cograph}
}

\newtheorem{theorem}{Theorem}
\newtheorem{lemma}[theorem]{Lemma}
\newtheorem{corollary}[theorem]{Corollary}
\theoremstyle{definition}
\newtheorem*{problem}{Problem 1}
\theoremstyle{remark}
\newtheorem*{remark}{Remark}

\newcommand{\Hstar}{H^{*}}
\newcommand{\Pfour}{P_{4}}

\title[A $P_4$-sparse minimal $(2,3)$-polar obstruction that is not a cograph]
{A nine-vertex $P_4$-sparse minimal $(2,3)$-polar obstruction\\ that is not a cograph}

\date{}

\subjclass[2020]{05C75, 05C17, 05C70}
\keywords{polarity, $(s,k)$-polar partition, minimal obstruction, $P_4$-sparse graph,
cograph, spider, counterexample}

\begin{document}

\begin{abstract}
Contreras-Mendoza and Hern\'andez-Cruz asked, as Problem~1 of their paper on minimal
obstructions for polarity in generalizations of cographs, whether every $P_4$-sparse minimal
$(s,k)$-polar obstruction is a cograph. We answer this in the negative at $(s,k)=(2,3)$. The
witness is the $9$-vertex, $15$-edge thin spider $\Hstar$ with head $2K_2+K_1$: it is
$P_4$-sparse, its unique induced $P_4$ shows that it is not a cograph, it is not
$(2,3)$-polar, and all nine of its one-vertex deletions are. Complementation gives a second
counterexample at $(s,k)=(3,2)$. The graph is exhibited by its edge list, and the proof is a
short case analysis together with nine explicit polar partitions, so no computation is needed
to check it.
\end{abstract}

\maketitle

\section{Introduction}

All graphs are finite, simple and undirected. A \emph{$k$-cluster} is a disjoint union of at
most $k$ cliques. For positive integers $s$ and $k$, a graph $G$ is \emph{$(s,k)$-polar} if
$V(G)$ has a partition $(A,B)$ such that $G[A]$ is a complete multipartite graph with at most
$s$ parts and $G[B]$ is a $k$-cluster; such a pair $(A,B)$ is an \emph{$(s,k)$-polar
partition}. Either side may be empty. Since $(1,1)$-polar is exactly split, polarity
generalises the F\"oldes--Hammer notion \cite{FoldesHammer}. The property is hereditary, so it
is characterised by its \emph{minimal obstructions}: graphs that are not $(s,k)$-polar although
every proper induced subgraph is. Consequently a graph is a minimal $(s,k)$-polar obstruction
as soon as it is not $(s,k)$-polar and all of its one-vertex deletions are.

A \emph{cograph} is a $P_4$-free graph, and a graph is \emph{$P_4$-sparse} if no five of its
vertices induce more than one $P_4$; every cograph is $P_4$-sparse. Contreras-Mendoza and
Hern\'andez-Cruz determined the minimal $(s,k)$-polar obstructions in these classes when
$\min\{s,k\}\le 1$ and posed the following question \cite[Problem~1]{CMHC24}.

\begin{problem}
For any positive integers $s$ and $k$, is every $P_4$-sparse minimal $(s,k)$-polar obstruction
a cograph?
\end{problem}

The problem carries this number both in the version of record of \cite{CMHC24} and in
version~2 of its arXiv e-print, and it is quoted verbatim above.

Two further facts from \cite{CMHC24} locate the question. First, a $P_4$-sparse minimal
$(s,k)$-polar obstruction has at most $(s+1)(k+1)$ vertices \cite[Theorem~57]{CMHC24}, a bound
the authors credit to Hannebauer's thesis \cite{Hannebauer}. Second, the authors' sequel proves
that there are exactly $50$ $P_4$-sparse minimal $2$-polar obstructions and that each of them
is a cograph \cite{CMHC22polar}; here $2$-polar means $(2,2)$-polar, and the minimal $2$-polar
obstructions inside the smaller class of cographs had been determined earlier
\cite{HellHCLS}. Together with the settled range $\min\{s,k\}\le1$, this leaves
$(s,k)=(2,3)$, where the bound reads $12$, as the smallest cell not covered by those results.
That is the cell we settle.

\begin{theorem}\label{thm:main}
Let $\Hstar$ be the graph on
$V=\{s_1,s_2,k_1,k_2,a_1,a_2,b_1,b_2,x\}$ whose edge set consists of exactly the following
fifteen pairs and no others:
\begin{center}
$s_1k_1$,\quad $s_2k_2$,\quad $k_1k_2$,\\
$k_1a_1$,\; $k_1a_2$,\; $k_1b_1$,\; $k_1b_2$,\; $k_1x$,\\
$k_2a_1$,\; $k_2a_2$,\; $k_2b_1$,\; $k_2b_2$,\; $k_2x$,\\
$a_1a_2$,\quad $b_1b_2$.
\end{center}
Then $\Hstar$ is $P_4$-sparse, it is not a cograph, and it is a minimal $(2,3)$-polar
obstruction. In particular, the answer to Problem~1 is \emph{no}.
\end{theorem}

Writing $S=\{s_1,s_2\}$, $K=\{k_1,k_2\}$ and $R=\{a_1,a_2,b_1,b_2,x\}$, the graph $\Hstar$ is
the spider \cite{JamisonOlariu} with independent set $S$, clique $K$, the two legs $s_1k_1$
and $s_2k_2$, and head
$R$, where $K$ is complete and $S$ is anticomplete to $R$ and $\Hstar[R]=2K_2+K_1$. Its degree
sequence is $(1,1,7,7,3,3,3,3,2)$, summing to $30=2\cdot15$, and the remaining $21$ of the
$\binom92=36$ pairs are non-edges. In \texttt{graph6} format, with the vertices in the order
$s_1,s_2,k_1,k_2,a_1,a_2,b_1,b_2,x$, the graph is the seven-character string
\begin{center}
\verb|HRKx`co|
\end{center}
and the same order presents the fifteen edges above as the $0$-indexed list
\begin{center}
\texttt{02 13 23 24 25 26 27 28 34 35 36 37 38 45 67}.
\end{center}
The graph is also $2K_2$-split: $V$ is the disjoint union of the clique $K$, the independent
set $\{s_1,s_2,x\}$, and the set $\{a_1,a_2,b_1,b_2\}$, which induces $2K_2$.

\begin{corollary}\label{cor:dual}
The complement $\overline{\Hstar}$ is a $P_4$-sparse minimal $(3,2)$-polar obstruction that is
not a cograph. Hence the answer to Problem~1 is also \emph{no} at $(s,k)=(3,2)$.
\end{corollary}

\begin{proof}
A graph $G$ is $(s,k)$-polar if and only if $\overline G$ is $(k,s)$-polar: complementation
exchanges ``complete multipartite with at most $s$ parts'' and ``disjoint union of at most $s$
cliques'', and it commutes with taking induced subgraphs, so it also exchanges minimal
$(s,k)$-polar obstructions with minimal $(k,s)$-polar obstructions. Complementation fixes the
set of induced $P_4$'s of a graph, because $\overline{P_4}\cong P_4$, so it preserves both
$P_4$-sparseness and the property of not being a cograph. Now apply
Theorem~\ref{thm:main}.
\end{proof}

\begin{remark}
What is refuted is Problem~1 as a statement quantified over all pairs $(s,k)$. The settled
range $\min\{s,k\}\le1$ and the case $(2,2)$ are theorems of \cite{CMHC24} and
\cite{CMHC22polar} and remain so, and nothing here bears on any other cell with
$\min\{s,k\}\ge2$, or on the classification, or even the finiteness, of the $P_4$-sparse
minimal $(2,3)$-polar obstructions; in particular we do not claim that $9$ is the least order
of a non-cograph example. Finally, $\Hstar$ is not among the $50$ obstructions of
\cite{CMHC22polar}, since $\Hstar-x$ is not $(2,2)$-polar and therefore $\Hstar$ is not a
minimal $(2,2)$-polar obstruction at all.
\end{remark}

\section{Proof of Theorem~\ref{thm:main}}

Throughout, $S$, $K$ and $R$ are as above. We use two elementary observations repeatedly:
every vertex of $K$ is adjacent to every vertex of $R$ and to the other vertex of $K$; and
$s_i$ has degree~$1$, its unique neighbour being $k_i$.

\medskip
\noindent\emph{Step 1: the only induced $\Pfour$ of $\Hstar$ is $s_1k_1k_2s_2$; hence $\Hstar$
is $P_4$-sparse and is not a cograph.}
The set $K\cup R$ induces the join of $K$ with $\Hstar[R]=2K_2+K_1$; a join of two cographs is
a cograph, so no induced $P_4$ lies inside $K\cup R$ and every induced $P_4$ of $\Hstar$ meets
$S$. Suppose $Q$ induces a $P_4$ and $s_1\in Q$. As $s_1$ has degree $1$, it is an endpoint of
that path and $k_1\in Q$. The third vertex of the path is a neighbour of $k_1$ other than
$s_1$, and the fourth is adjacent to the third but not to $k_1$; the only vertex of $\Hstar$
other than $s_1$ that is not adjacent to $k_1$ is $s_2$, so the fourth vertex is $s_2$ and the
third is its unique neighbour $k_2$. Thus $Q=\{s_1,k_1,k_2,s_2\}$, and symmetrically if
$s_2\in Q$. Conversely $s_1k_1k_2s_2$ is an induced path, since $s_1s_2$, $s_1k_2$ and
$s_2k_1$ are non-edges. So $\Hstar$ has exactly one induced $P_4$: no five vertices can
contain two of them, so $\Hstar$ is $P_4$-sparse, and $\Hstar$ is not $P_4$-free, so it is not
a cograph.

\medskip
For the remaining two steps, note that a complete multipartite graph with at most two parts is
exactly a complete bipartite graph $K_{p,q}$ with $p+q\ge1$, $p,q \ge 0$; in particular it is
triangle-free. We record the only property of it we need.

\begin{lemma}\label{lem:cb}
Let $G[A]$ be complete multipartite with at most two parts and let $uv$ be an edge with
$u,v\in A$. Then every $w\in A$ is adjacent to $u$ or to $v$.
\end{lemma}

\begin{proof}
Parts are independent, so $u$ and $v$ lie in different parts, and there are at most two parts,
so $w$ lies in the part of $u$ or in the part of $v$. Distinct parts are completely joined, so
in the first case $w$ is adjacent to $v$ and in the second to $u$.
\end{proof}

\begin{lemma}\label{lem:B}
Let $(A,B)$ be a partition of $V(\Hstar)$ with $\Hstar[B]$ a $3$-cluster. If $B\cap K \ne
\emptyset$, then $B\cap R$ is a clique, and consequently $\{a_1,a_2\}\subseteq A$ or
$\{b_1,b_2\}\subseteq A$.
\end{lemma}

\begin{proof}
Let $k\in B\cap K$. As $k$ is adjacent to every vertex of $R$, the set $(B\cap R)\cup\{k\}$
lies in a single connected component of $\Hstar[B]$, and components of a cluster are cliques;
hence $B \cap R$ is a clique. The cliques of $\Hstar[R]=2K_2+K_1$ are the subsets of
$\{a_1,a_2\}$, of $\{b_1,b_2\}$ and of $\{x\}$, so $B\cap R$ misses one of the two pairs
$\{a_1,a_2\}$, $\{b_1,b_2\}$ entirely, and that pair lies in $A$.
\end{proof}

\medskip
\noindent\emph{Step 2: $\Hstar$ is not $(2,3)$-polar.}
Suppose $(A,B)$ were a $(2,3)$-polar partition. We distinguish two cases.

\emph{Case $B\cap K\ne\emptyset$.} By Lemma~\ref{lem:B} the set $A$ contains one of the edges
$a_1a_2$, $b_1b_2$; call it $uv$. If the other vertex $k'$ of $K$ lies in $A$, then $k'$ is
adjacent to both $u$ and $v$, so $\{k',u,v\}$ is a triangle in $A$, contradicting the fact that
$\Hstar[A]$ is complete bipartite. So $K\subseteq B$. Then $s_1\in A$: for if $s_1\in B$, its
unique neighbour $k_1$ lies in $B$ as well, so $s_1$ and $k_1$ are in one component of
$\Hstar[B]$, a clique that also contains $k_2$ because $k_1k_2$ is an edge, while $s_1k_2$ is a
non-edge --- impossible. But $s_1$ is adjacent to neither $u$ nor $v$, since $S$ is
anticomplete to $R$, and this contradicts Lemma~\ref{lem:cb} applied to the edge $uv$ of $A$.

\emph{Case $K\subseteq A$.} No vertex $r\in R$ can lie in $A$, since $\{k_1,k_2,r\}$ would be a
triangle; hence $R\subseteq B$. The three components of $\Hstar[R]$ are pairwise anticomplete
and $K\cap B=\emptyset$, so they are already three distinct components of $\Hstar[B]$. If some
$s_i$ lay in $B$, it would be a fourth component, its unique neighbour $k_i$ being in $A$; then
$\Hstar[B]$ would have at least four components and could not be a $3$-cluster. So
$S\subseteq A$ and $A=\{s_1,s_2,k_1,k_2\}$, which induces the path $s_1k_1k_2s_2$. Apply
Lemma~\ref{lem:cb} to the edge $s_1k_1$ of $A$ and to $w=s_2$: the vertex $s_2$ would have to be
adjacent to $s_1$ or to $k_1$, and it is adjacent to neither.

Both cases are impossible, so no $(2,3)$-polar partition of $\Hstar$ exists.

\medskip
\noindent\emph{Step 3: every one-vertex deletion of $\Hstar$ is $(2,3)$-polar.}
The nine partitions below are $(2,3)$-polar partitions of the nine graphs $\Hstar-v$. For each
we list $A$ split into its parts, and $B$ split into its clique components, so that each line
can be checked against the definition by inspection: the parts are independent sets pairwise
completely joined, the components are cliques pairwise anticomplete, and $A\cup B=V\setminus
\{v\}$.

\medskip
\begin{center}
\renewcommand{\arraystretch}{1.15}
\begin{tabular}{llll}
$v$ & parts of $A$ & clique components of $B$ & \\
\hline
$s_1$ & $\{k_2\}\mid\{k_1,s_2\}$ & $\{a_1,a_2\}\mid\{b_1,b_2\}\mid\{x\}$ & \\
$s_2$ & $\{k_1\}\mid\{k_2,s_1\}$ & $\{a_1,a_2\}\mid\{b_1,b_2\}\mid\{x\}$ & \\
$k_1$ & $\{k_2\}\mid\{s_2,x\}$   & $\{a_1,a_2\}\mid\{b_1,b_2\}\mid\{s_1\}$ & \\
$k_2$ & $\{k_1\}\mid\{s_1,x\}$   & $\{a_1,a_2\}\mid\{b_1,b_2\}\mid\{s_2\}$ & \\
$x$   & $\{k_1\}\mid\{k_2,s_1\}$ & $\{a_1,a_2\}\mid\{b_1,b_2\}\mid\{s_2\}$ & \\
$a_1$ & $\{s_1,s_2,a_2,x\}$      & $\{k_1,k_2,b_1,b_2\}$ & \\
$a_2$ & $\{s_1,s_2,a_1,x\}$      & $\{k_1,k_2,b_1,b_2\}$ & \\
$b_1$ & $\{s_1,s_2,b_2,x\}$      & $\{k_1,k_2,a_1,a_2\}$ & \\
$b_2$ & $\{s_1,s_2,b_1,x\}$      & $\{k_1,k_2,a_1,a_2\}$ & \\
\end{tabular}
\end{center}

\medskip
For the first five rows, $A$ has two parts and $B$ is a $3$-cluster; for the last four, $A$ is
independent (one part) and $B$ is the clique $K_4$ induced by $K$ together with the surviving
pair of the head. Two rows deserve a word. In the row $v=k_1$ the vertex $s_1$ becomes isolated
because its only neighbour has been deleted, which is why it can be a clique component of $B$
on its own; the pair $\{s_2,x\}$ is independent and both of its vertices are adjacent to $k_2$,
so $\Hstar[A]$ is $K_{1,2}$. In the row $v=a_1$ the set $\{s_1,s_2,a_2,x\}$ is independent,
because $S$ is anticomplete to $R$ and $a_2$'s only neighbour inside $R$ was $a_1$; and
$\{k_1,k_2,b_1,b_2\}$ is a $K_4$ since $K$ is a clique complete to $R$ and $b_1b_2$ is an edge.

By Step~2 and Step~3, $\Hstar$ is a minimal $(2,3)$-polar obstruction, and by Step~1 it is
$P_4$-sparse and not a cograph. \qed

\section{Verification}

Everything asserted about $\Hstar$ in Theorem~\ref{thm:main} is decided by the hand argument of
Section~2 from the edge list printed in Section~1, and a reader who wants to check it needs no
computer. The claims were nevertheless re-derived independently, by a program written from the
printed strings alone and using exact bitmask and integer arithmetic with no third-party
library: it decodes the \texttt{graph6} string, checks that the decoded edge set is
\emph{equal} to the printed labelled edge list, and then recomputes the degree sequence, the
spider structure, the single induced $P_4$ over all $\binom94=126$ four-subsets, the
$P_4$-sparseness over all $\binom95=126$ five-subsets, the failure of $(2,3)$-polarity by
exhausting all $2^9=512$ bipartitions, the nine deletion certificates of Step~3 one at a time
against the definition, and independently of them a $(2,3)$-polar partition of each deletion
found by exhaustion; it also recomputes Corollary~\ref{cor:dual} on the complement. All $71$ of
its checks passed. The program \texttt{verify.py} and the transcript
\texttt{verify.output.txt} of the run accompany this note; they are not needed in order to
check the theorem.

\medskip
\noindent\textbf{Relation to the literature.}
The nearest published work is the same authors' characterisation of polarity on $H$-split
graphs \cite{CMHCsplit}: $\Hstar$ is a $2K_2$-split graph, it satisfies the order bound
$s+2k+2=10$ proved there, and that characterisation confirms clause by clause both the failure
of $(2,3)$-polarity and the nine certificates above; it is not applied there to $P_4$-sparse
graphs or to cographs. We make no priority claim. In particular we have not inspected
\cite{Hannebauer}, the 2010 thesis on matrix colorings of $P_4$-sparse graphs from which the
order bound $(s+1)(k+1)$ is quoted in \cite{CMHC24}.

\begin{thebibliography}{9}

\bibitem{CMHC24}
F.~E. Contreras-Mendoza and C.~Hern\'andez-Cruz,
\emph{Minimal obstructions for polarity, monopolarity, unipolarity and $(s,1)$-polarity in
generalizations of cographs},
Graphs Combin. \textbf{40} (2024), no.~3, Paper No.~53.
\href{https://doi.org/10.1007/s00373-024-02784-7}{doi:10.1007/s00373-024-02784-7};
\href{https://arxiv.org/abs/2203.04953}{arXiv:2203.04953}.

\bibitem{CMHC22polar}
F.~E. Contreras-Mendoza and C.~Hern\'andez-Cruz,
\emph{$2$-polarity and algorithmic aspects of polarity variants on cograph superclasses},
Discrete Math. Theor. Comput. Sci.
\href{https://doi.org/10.46298/dmtcs.11479}{doi:10.46298/dmtcs.11479};
\href{https://arxiv.org/abs/2210.02497}{arXiv:2210.02497}.

\bibitem{CMHCsplit}
F.~E. Contreras-Mendoza and C.~Hern\'andez-Cruz,
\emph{Polarity on $H$-split graphs},
Discrete Appl. Math. \textbf{384} (2026), 229--242.
\href{https://doi.org/10.1016/j.dam.2025.12.065}{doi:10.1016/j.dam.2025.12.065};
\href{https://arxiv.org/abs/2303.17055}{arXiv:2303.17055}.
An earlier version appeared in Procedia Comput. Sci. \textbf{223} (2023), 184--192.

\bibitem{FoldesHammer}
S.~F\"oldes and P.~L. Hammer,
\emph{Split graphs},
in: Proceedings of the Eighth Southeastern Conference on Combinatorics, Graph Theory and
Computing, Congress. Numer. \textbf{XIX}, 1977, pp.~311--315.

\bibitem{Hannebauer}
C.~Hannebauer,
\emph{Matrix colorings of $P_4$-sparse graphs},
Master's thesis, FernUniversit\"at in Hagen, 2010.
Cited as the source of the bound $(s+1)(k+1)$ in \cite[Theorem~57]{CMHC24}.

\bibitem{HellHCLS}
P.~Hell, C.~Hern\'andez-Cruz and A.~Linhares-Sales,
\emph{Minimal obstructions to $2$-polar cographs},
Discrete Appl. Math. \textbf{261} (2019).

\bibitem{JamisonOlariu}
B.~Jamison and S.~Olariu,
\emph{A tree representation for $P_4$-sparse graphs},
Discrete Appl. Math. \textbf{35} (1992), no.~2, 115--129.
\href{https://doi.org/10.1016/0166-218X(92)90036-B}{doi:10.1016/0166-218X(92)90036-B}.

\end{thebibliography}

\end{document}
